\section{Instantiations}\label{sec:prelims-inst}

\begin{outline}
  \item what this section gives: the groups, assumptions and concrete schemes the rest builds on
  \item what is included here: the standard instantiation of a primitive from this chapter (Pedersen, ElGamal), or prior work used in more than one chapter ($\AGHO$, BBS+)
  \item schemes used in one chapter only, variants modified for one chapter, and our own constructions are in the chapter that uses them
\end{outline}

\subsection{Groups and assumptions}\label{sec:prelims-groups}

\begin{outline}
  \item $\Zp$ for prime $p$: integers mod $p$, elements called field elements; $\Zp^* = \Zp \setminus \{0\}$
  \item cyclic groups of prime order $p$, written multiplicatively; their elements are group elements
  \item discrete logarithm: for generator $\gen$ and $a = \gen^x$, $x \in \Zp$ is the discrete logarithm of $a$ relative to $\gen$
\end{outline}

Most of our schemes use three such groups, $\group{1}$, $\group{2}$, and
$\group{T}$ of prime order $p$, with generators $\gen \in \group{1}$ and $\ggen
\in \group{2}$. The groups are equipped with an efficiently computable pairing
$\pairing \colon \group{1} \times \group{2} \rightarrow \group{T}$. For all $a
\in \group{1}$, $\widetilde{b} \in \group{2}$, and $x,y \in \Zp$, bilinearity
means that the following holds:
%
\begin{equation*}
%
  \pairing(a^x,\widetilde{b}^y) = \pairing(a,\widetilde{b})^{xy}.
%
\end{equation*}
%
The pairing is non-degenerate, so $\pairing(\gen,\ggen)$ generates $\group{T}$.
We call $\group{1}$ and $\group{2}$ the source groups and $\group{T}$ the target
group. We work in the asymmetric, or type III, setting, in which $\group{1} \neq
\group{2}$ and no efficiently computable homomorphism between them is known.
Elements of $\group{2}$ carry a tilde, and field elements are usually written
with Greek letters.

Bilinearity allows relations between exponents in the source groups to be
checked through equalities in $\group{T}$. The equality
$\pairing(a,\ggen) = \pairing(\gen,\widetilde{b})$ shows that $a$ and
$\widetilde{b}$ have the same discrete logarithm relative to their generators,
without that discrete logarithm being recovered. A pairing product equation is
an equation of the form $\prod_j \pairing(a_j, \widetilde{b}_j) = t$ for $t \in
\group{T}$. Our constructions use such equations to verify signatures and
algebraic relations across $\group{1}$ and $\group{2}$.

\begin{outline}
  \item we rely on hardness of discrete logarithms (\Cref{asm:dl}) and on DDH (\Cref{asm:ddh}), cite DifHel76, Boneh98
  \item assumptions specific to one chapter are given where they are used
\end{outline}

\begin{assumptionbox}{Discrete logarithm (DL)}{dl}

  The DL assumption holds in a group of prime order $p$ with generator $\gen$
if, for $x \sample \Zp$, no PPT adversary given $\gen^x$ outputs $x$ with
non-negligible probability.

\end{assumptionbox}

\begin{assumptionbox}{Decisional Diffie-Hellman (DDH)}{ddh}

  The DDH assumption holds in a group of prime order $p$ with generator $\gen$
if $(\gen^x, \gen^y, \gen^{xy})$ and $(\gen^x, \gen^y, \gen^z)$, for $x, y, z
\sample \Zp$, are computationally indistinguishable.

\end{assumptionbox}

The symmetric external Diffie-Hellman (SXDH) assumption holds in a bilinear group
if DDH holds in both $\group{1}$ and $\group{2}$.

\begin{outline}
  \item DDH cannot hold in a symmetric bilinear group ($\group{1} = \group{2}$), since the pairing decides it; so SXDH needs type III
  \item some schemes we use are only proven in the generic group model (GGM), cite EC:Shoup97, IMA:Maurer05
  \item GGM: adversary sees only opaque handles to group elements, computes only through oracles for the group operations and the pairing
  \item a GGM proof rules out attacks that ignore the representation of the group; weaker than a reduction to a stated assumption
  \item $q$-type assumptions: statement grows with a parameter $q$ (typically the number of signing queries); stronger than fixed-size ones like DDH
\end{outline}

\subsection{Sigma protocols and the Fiat-Shamir transform}\label{sec:prelims-sigma}

\begin{outline}
  \item most proof systems here are Sigma protocols made non-interactive
  \item Sigma protocol for $\rel$: three moves, prover sends $a$, verifier sends uniform challenge $c$, prover sends response $z$, verifier checks against $(\inst, a, c)$
  \item completeness: as for proof systems
  \item special soundness: efficient extractor gets a witness from two accepting transcripts $(a, c, z)$, $(a, c', z')$ with $c \neq c'$
  \item special honest-verifier zero knowledge: efficient simulator, given $\inst$ and $c$, outputs $(a, z)$ distributed as an honest run with challenge $c$ (cite C:CraDamSch94)
\end{outline}

\begin{outline}
  \item example: Schnorr's protocol for knowledge of $x$ with $h = \gen^x$ (cite JC:Schnorr91)
  \item prover sends $a = \gen^r$ for $r \sample \Zp$, gets $c \sample \Zp$, returns $z = r + cx$; verifier checks $\gen^z = a h^c$
  \item same protocol proves a representation $h = \prod_i \gen_i^{x_i}$, one response per exponent; non-interactive form = Schnorr proof
  \item several equations over the same witness: run in parallel with one shared challenge, so they hold together (e.g.\ two group elements with the same discrete logarithm)
\end{outline}

\begin{outline}
  \item Fiat-Shamir: prover computes $c = \hash(\inst, a, \lbl)$ itself, $\hash$ a random oracle, $\lbl$ a label (cite C:FiaSha86)
  \item quasi-unique responses: for given $a$ and $c$, no efficient prover finds two accepting responses; unique responses (stronger): response fixed by a verification equation, as in Schnorr
  \item result: whole statement hashed + special soundness + SHVZK + quasi-unique responses gives a simulation extractable labelled proof system in the ROM (cite INDOCRYPT:FKMV12)
  \item why hash the whole statement: a part left out can be chosen after seeing the challenge
  \item ZK simulator programs the random oracle at $(\inst, a, \lbl)$
  \item extractor rewinds: reruns $\adv$ from the query that fixed the challenge, answers with a fresh challenge, applies special soundness; forking lemma bounds success (cite JC:PoiSte00)
  \item extraction from a single run without rewinding = straight-line/online, what UC needs (cite C:Fischlin05; links back to \Cref{sec:prelims-security})
\end{outline}

\subsection{The Pedersen commitment scheme}

\begin{outline}
  \item \Cref{prim:pedersen}: commits to several field elements at once with one group element (cite C:Pedersen91)
\end{outline}

\begin{primitivebox}{The Pedersen commitment scheme}{pedersen}

  \begin{description}

    \item[$\ck \sample \ComGen(\secparam, n)$ :] On input the security parameter
and message vector length $n$, sample $\gen_1,\ldots,\gen_n, \hgen \sample
\group{1}$ and return $\ck = (\gen_1,\ldots,\gen_n,\hgen)$.

    \item[$\com \leftarrow \ComCommit(\ck,(m_1,\ldots,m_n),\openinfo)$ :] On
input commitment key $\ck$, messages $(m_1,\ldots,m_n) \in \Zp^n$, and opening
$\openinfo \in \Zp$, return $\com = \prod_{i=1}^{n} \gen_i^{m_i}\cdot
\hgen^{\openinfo}$.

    \item[$(m_1,\ldots,m_n) \leftarrow \ComOpen(\ck,\com,\openinfo)$ :] On
input $\ck$, commitment $\com$, and opening $\openinfo = (m_1,\ldots,m_n,
\openinfo')$, return $(m_1,\ldots,m_n)$ if $\com = \prod_{i=1}^{n}
\gen_i^{m_i}\cdot\hgen^{\openinfo'}$, and $\bot$ otherwise.

  \end{description}

\end{primitivebox}

The Pedersen commitment scheme is perfectly hiding, and computationally binding
under the DL assumption~\cite{C:Pedersen91}.

\subsection{The ElGamal encryption scheme}

\begin{outline}
  \item \Cref{prim:elgamal}: encrypts group elements (cite C:ElGamal84)
\end{outline}

\begin{primitivebox}{The ElGamal encryption scheme}{elgamal}

  \begin{description}

    \item[$(\pk, \sk) \sample \PKEgen(\secparam)$ :] Sample $\mu \sample \Zp$,
and return $\sk = \mu$ and $\pk = \gen^{\mu}$.

    \item[$\ciph \sample \PKEencrypt(\pk, M)$ :] On input $M \in \group{1}$,
sample $\rho \sample \Zp$, and return $\ciph = (\gen^{\rho}, M \pk^{\rho})$.

    \item[$M \leftarrow \PKEdecrypt(\sk, \ciph)$ :] On input $\ciph = (f_1, f_2)$,
return $f_2 / f_1^{\mu}$.

  \end{description}

\end{primitivebox}

\begin{outline}
  \item ElGamal is IND-CPA under DDH in $\group{1}$ (cite C:ElGamal84)
  \item messages are group elements: a field element $m$ is encrypted as $\gen^m$, and decryption gives back $\gen^m$, not $m$
\end{outline}

\subsection{Structure-preserving signature schemes}\label{sec:prelims-sps}

\begin{outline}
  \item SPS schemes: signatures where messages, public keys and signatures are vectors of group elements and verification only checks pairing product equations (cite C:AFGHO10, C:AGHO11)
  \item why we want them: nothing is hashed, so they combine with efficient zero-knowledge proofs over the same groups; users prove knowledge of signatures on committed or hidden group elements
\end{outline}

An SPS scheme $\SPS$ over the bilinear groups above consists of the following
algorithms:

\begin{description}

  \item[$(\pk, \sk) \sample \SPSgen(\secparam, \ell)$ :] The key generation
algorithm, on input the security parameter and a vector length $\ell$, outputs
public verification key $\pk$, which consists of group elements, and secret
signing key $\sk$.

  \item[$\SPSsig \sample \SPSsign(\sk, \vec{\msg})$ :] The signing algorithm
signs message $\vec{\msg} \in \group{1}^\ell$ using secret key $\sk$, and outputs
a signature $\SPSsig$ consisting of elements of $\group{1}$ and $\group{2}$.

  \item[$\{0,1\} \leftarrow \SPSverify(\pk, \vec{\msg}, \SPSsig)$ :]
The verification algorithm evaluates pairing product equations in $\pk$,
$\vec{\msg}$ and $\SPSsig$, and outputs $0$ or $1$.

\end{description}
%
Some schemes are rerandomisable: they additionally provide an algorithm
$\SPSsig' \sample \SPSrandomise(\SPSsig)$, which, without the secret key,
turns a signature into a new signature on the same message, distributed as a
freshly generated one.

\begin{outline}
  \item a rerandomisable scheme cannot be strongly unforgeable: rerandomising is itself a strong forgery
\end{outline}

Correctness and EUF-CMA security are defined as for digital signature schemes,
with messages restricted to $\group{1}^\ell$.

\subsection{The $\AGHO$ SPS scheme}

\begin{outline}
  \item Abe, Groth, Haralambiev and Ohkubo: SPS signatures in type III groups have at least three group elements; they give schemes meeting this (cite C:AGHO11)
  \item their main scheme: strongly unforgeable in the GGM; a six-element variant: strongly unforgeable under a non-interactive $q$-type assumption
  \item we use their rerandomisable three-element scheme, $\group{1}$ and $\group{2}$ swapped so messages are in $\group{1}$, and call it $\AGHO$
  \item secret key is only field elements; the message enters signing only through the terms $\msg_i^{-\mu_i}$ (matters for the threshold version in \Cref{chap:tsps})
  \item point to \Cref{prim:agho}
\end{outline}

\begin{primitivebox}{The $\AGHO$ SPS scheme}{agho}

  \begin{description}

    \item[$(\pk, \sk) \sample \SPSgen(\secparam, \ell)$ :] Sample $\mu_1,
\ldots, \mu_\ell, \tau \sample \Zp^*$, and output $\sk = (\mu_1, \ldots, \mu_\ell,
\tau)$ and $\pk = (\widetilde{u}_1, \ldots, \widetilde{u}_\ell, v)$, where
$\widetilde{u}_i = \ggen^{\mu_i}$ and $v = \gen^{\tau}$.

    \item[$\SPSsig \sample \SPSsign(\sk, \vec{\msg})$ :] On input $\vec{\msg} =
(\msg_1, \ldots, \msg_\ell) \in \group{1}^\ell$, sample $\omega \sample \Zp^*$
and output
      %
      \begin{equation*}
      %
	\SPSsig = (\widetilde{r}, \widetilde{s}, z) = \left(\ggen^{\omega}, \
\widetilde{r}^{\,\tau}, \ \left(\gen \prod_{i=1}^{\ell}
\msg_i^{-\mu_i}\right)^{1/\omega}\right).
      %
      \end{equation*}

    \item[$\SPSsig' \sample \SPSrandomise(\SPSsig)$ :] On input $\SPSsig =
(\widetilde{r}, \widetilde{s}, z)$, sample $\omega' \sample \Zp^*$ and output
$(\widetilde{r}^{\,\omega'}, \widetilde{s}^{\,\omega'}, z^{1/\omega'})$.

    \item[$\{0,1\} \leftarrow \SPSverify(\pk, \vec{\msg}, \SPSsig)$ :] On input
$\SPSsig = (\widetilde{r}, \widetilde{s}, z)$, output $1$ if and only if
      %
      \begin{equation*}
      %
	\pairing(v, \widetilde{r}) = \pairing(\gen, \widetilde{s}) \ \wedge\
\pairing(z, \widetilde{r}) \prod_{i=1}^{\ell} \pairing(\msg_i, \widetilde{u}_i)
= \pairing(\gen, \ggen).
      %
      \end{equation*}

  \end{description}

\end{primitivebox}

\begin{outline}
  \item $\AGHO$ is EUF-CMA secure in the GGM (cite C:AGHO11)
  \item being rerandomisable, it is not strongly unforgeable
\end{outline}

\subsection{The BBS+ signature scheme}\label{sec:prelims-bbs}

The BBS+ signature scheme of \Cref{prim:bbs} signs several messages at
once~\cite{C:BonBoySha04,SCN:AuSusMu06,EPRINT:CamDriLeh16}.

\begin{primitivebox}{The BBS+ signature scheme}{bbs}

  \begin{description}

    \item[$(\sk,\pk) \sample \BBSgen(\secparam, n, (\gen_0,\ldots,\gen_{n+1},
\ggen))$ :] On input the security parameter, vector length $n$, and generators
$\gen_0,\gen_1, \ldots, \gen_{n+1} \in \group{1}$, $\ggen \in \group{2}$, do the
following:

      \begin{enumerate}

	\item Sample $z \sample \Zp$.

	\item Return $\sk = z$, $\pk = (\gen_0, \ldots, \gen_{n+1}, \ggen,
\ggen_Z = \ggen^z)$.

      \end{enumerate}

      The generators are an input rather than an output. Standard presentations
have the signer sample them, which is sound when the signer is the only party
relying on them. Where the same generators also key a commitment whose binding
must hold against the signer, they are instead fixed by hashing to $\group{1}$
from a public seed, so that no party knows a discrete logarithm relation among
them.

    \item[$\signature \sample \BBSsign(\sk,(\msg_1, \ldots, \msg_n))$ :] On
input secret key $\sk = z$ and message $(\msg_1, \ldots, \msg_n) \in
\group{1}^n$, do:

      \begin{enumerate}

	\item Sample $e,r \sample \Zp$.

	\item Return $\signature = (e, r, E)$ with $E = (\gen_0 \cdot
\prod_{i=1}^{n} \msg_i \cdot \gen_{n+1}^{r})^{1/(z+e)}$.

      \end{enumerate}

    \item[$\{0,1\} \leftarrow \BBSverify(\pk,(\msg_1, \ldots, \msg_n),
\signature)$ :] On input public key $\pk$, message $(\msg_1, \ldots, \msg_n)$,
and signature $\signature = (e,r,E)$, do:

      \begin{enumerate}

	\item Return $\pairing(E,\ggen_Z\cdot\ggen^{e}) \checkcheck
\pairing(\gen_0 \cdot \prod_{i=1}^{n} \msg_i \cdot \gen_{n+1}^{r}, \ggen)$.

      \end{enumerate}

  \end{description}

\end{primitivebox}

The BBS+ signature scheme is existentially unforgeable under the $q$-strong
Diffie-Hellman ($q$-SDH) assumption of \Cref{asm:qsdh-bbs}, a $q$-type
assumption~\cite{EC:BonBoy04a}.

\begin{assumptionbox}{$q$-Strong Diffie-Hellman ($q$-SDH)}{qsdh-bbs}

  The $q$-SDH assumption holds in the type III bilinear group
$(\group{1},\group{2},\group{T},\pairing)$ if, given $(\gen_0, \gen_0^{\theta},
\ldots, \gen_0^{\theta^{q}}, \ggen, \ggen^{\theta})$ for $\gen_0 \in \group{1}$,
$\ggen \in \group{2}$ and uniformly random $\theta \sample \Zp$, no PPT
adversary can output a pair $\big(\varsigma,\,
\gen_0^{1/(\theta+\varsigma)}\big)$ for any $\varsigma \in \Zp \setminus
\{-\theta\}$ with non-negligible probability.

\end{assumptionbox}

This presentation signs group elements. To sign field elements $m_1, \ldots, m_n
\in \Zp$, the signer signs $M_i = \gen_i^{m_i}$. Signing the group elements
$M_i$ directly, rather than the field elements $m_i$ themselves, is useful
because computing $\prod_i M_i$ requires no knowledge of the discrete logarithms $m_i$,
so a signer can produce a valid signature on committed values without ever
learning them~\cite{SCN:AuSusMu06,EC:TesZhu23b}.

Unforgeability when signing group elements carries a condition that signing
field elements does not. An adversary able to obtain a signature on an arbitrary group
element could ask
for one on $M_i' = M_i \cdot \gen_i^{\delta}$, shifting the certified message by
a quantity it knows. The reduction to $q$-SDH therefore requires the signer to
sign only elements that come with a proof of knowledge of their representation
in the bases $\gen_i$~\cite{SCN:AuSusMu06,EC:TesZhu23b}.
